\chapter{Results and Comparative Analysis}
{\itshape This chapter presents the experimental evaluation of the classical and quantum-inspired K-Means segmentation frameworks. The performance is assessed across hardware metrics (latency, throughput, and GPU utilization) and mathematical clustering indicators (Within-Cluster Sum of Squares, Davies-Bouldin index, and Silhouette score). Visual output quality and temporal boundary stability are analyzed under varying lighting conditions.}

\section{System Throughput and Latency Scaling}
This section evaluates the execution speed of both clustering engines under different computational loads, modified by the number of clusters ($K$) and the data downsampling stride ($S$).

\subsection{Hyperparameter Latency Benchmarks}
To measure how the algorithms scale, tests were conducted on a static $640 \times 480$ video frame depicting a complex desk layout under constant studio lighting. To ensure a scientifically fair comparison and isolate the mathematical execution paths of both engines, all trials were conducted using the K-Means++ initialization algorithm. For each benchmark run, both engines were seeded with the exact same set of initial starting centroids and executed on the identical image frame. Table~\ref{tab:latency_benchmarks} compares the average core engine execution latency and the convergence iteration counts over 21 runs for three representative configurations.

\begin{table}[H]
\centering
\caption{Performance comparison of Classical vs. Quantum K-Means under different configurations.}
\label{tab:latency_benchmarks}
\begin{tabular}{ccccccc}
\hline
\multicolumn{2}{c}{\textbf{Configuration}} & \multicolumn{2}{c}{\textbf{Classical Engine}} & \multicolumn{2}{c}{\textbf{Quantum Engine}} \\
\textbf{K} & \textbf{Stride (S)} & \textbf{Latency (ms)} & \textbf{Iterations} & \textbf{Latency (ms)} & \textbf{Iterations} \\ \hline
3  & 8  & 1.8056  & 11.71  & 1.6107  & 12.38  \\
8  & 4  & 6.4545  & 41.81  & 6.1547  & 41.43  \\
40 & 1  & 26.6608 & 88.67  & 27.9470 & 88.62  \\ \hline
\end{tabular}
\end{table}

The latency data shows that the execution time scales non-linearly with both $K$ and $S$. In the low-density configuration ($K=3, S=8$), both engines execute in under $2\text{ ms}$. In the high-density configuration ($K=40, S=1$), which processes all $307,200$ pixels without downsampling, latency increases to approximately $26.6\text{ ms}$ for the classical engine and $27.9\text{ ms}$ for the quantum engine. Despite the mathematical complexity of the quantum-inspired state projection, the execution latency of both engines remains highly comparable across all modes.

\subsection{GPU Latency-Hiding and Hardware Intrinsics}
The small latency difference between the classical and quantum-inspired engines is explained by the architecture of modern NVIDIA GPUs:

\begin{enumerate}
    \item \textbf{Special Function Units (SFUs)}: The classical engine relies on standard Arithmetic Logic Units (ALUs) for basic Euclidean distance math (multiplications and additions). The quantum engine uses trigonometric phase projections, which require cosine calculations. The CUDA compiler (\texttt{nvcc}) compiles these operations into the hardware-accelerated intrinsic \texttt{\_\_cosf()}, which executes directly on the GPU's dedicated Special Function Units (SFUs) in a small number of clock cycles rather than using general ALU emulation.
    
    \item \textbf{Memory Bandwidth Bottleneck}: For both engines, the execution time is dominated by global memory access (reading pixel values and centroid locations) rather than raw arithmetic logic. Because the GPU can hide arithmetic latency by scheduling threads while waiting for memory fetches to complete (latency-hiding), the extra trigonometric calculations of the quantum engine have a minimal impact on the overall execution time.
    
    \item \textbf{Camera I/O Bottleneck and GPU Utilization}: During real-time execution on an NVIDIA GeForce RTX 5070 Ti, the average GPU utilization hovers around 40\%, occasionally spiking to 43\%, in the heaviest configuration ($K=40, S=1$). This behavior is caused by the camera's hardware frame rate, which is capped at 30 frames per second. Since the camera delivers a new frame every $33.3\text{ ms}$ and the engine requires only $26.6\text{ ms}$ (classical) or $27.9\text{ ms}$ (quantum) to process it, the GPU remains idle for the remaining portion of the frame interval. In the lighter configuration ($K=3, S=8$), the core execution time drops to under $2\text{ ms}$, meaning the GPU remains idle for over 90\% of the frame window. In this mode, the average GPU utilization stays around 30--35\%, which primarily represents the fixed overhead of rendering the ImGui interface and maintaining the OpenGL window context. This operational constraint justifies the use of offline metrics and overall throughput calculations to estimate the maximum performance boundaries of the hardware.
\end{enumerate}